\scanpage{253}
% archive sheet number: 252
\margin{\handwritten{Grothendieck}}% short separate horizontal impression below the inscription

\hfill n$^{\mathrm{o}}$ \handwritten{293}

\begin{center}
\uline{LE THEOREME PRINCIPAL DE Z A R I S K I}

I. - LEMMES PRELIMINAIRES
\end{center}

\uline{Lemme} 1 - \uline{Soient} $\mathrm r$ \uline{un domaine d'intégrité et} $\mathrm o$ \uline{un sous-anneau de} $\mathrm r$. \uline{On suppose qu'il existe des éléments} $t_1,..,t_n$ \uline{de} $\mathrm r$ $(n>0)$ \uline{qui sont algébriquement indépendants par rapport à} $\mathrm o$ \uline{et tels que} $\mathrm r$ \uline{soit entier sur} $\mathrm o[t_1,....,t_n]$. \uline{Soient} $\mathrm q$ \uline{un idéal premier maximal de} $\mathrm r$ \uline{et} $\mathrm p=\mathrm q\cap\mathrm o$. \uline{Il y a alors un idéal premier de} $\mathrm r$ \uline{contenu dans} $\mathrm q$, \uline{distinct de} $\mathrm q$ \uline{et contenant} $\mathrm p$.

Nous nous ramènerons d'abord au cas où $\mathrm o$ est intégralement clos. Soit $\mathrm o^*$ la clôture intégrale de $\mathrm o$, et soit $\mathrm r^*=\mathrm r[\mathrm o^*]$; les éléments $t_1,...,t_n$ sont algébriquement indépendants sur $\mathrm o^*$ et $\mathrm r^*$ est entier sur $\mathrm o^*[t_1,...,t_n]$. Soit $\mathrm q^*$ un idéal premier de $\mathrm r^*$ tel que $\mathrm q^*\cap\mathrm r=\mathrm q$; comme $\mathrm r^*$ est entier sur $\mathrm r$, $\mathrm q^*$ est un idéal premier maximal de $\mathrm r^*$. Soit $\mathrm p^*=\mathrm q^*\cap\mathrm o^*$. Supposons qu'il y ait un idéal premier $\mathrm n^*$ de $\mathrm r^*$ contenu dans $\mathrm q^*$, distinct de $\mathrm q^*$ et contenant $\mathrm p^*$. Soit $\mathrm n=\mathrm n^*\cap\mathrm r$; $\mathrm n$ est un idéal premier de $\mathrm r$; comme $\mathrm q^*\cap\mathrm r=\mathrm q$ on a $\mathrm n\subset\mathrm q$; il résulte du fait que $\mathrm n^*\ne\mathrm q^*$ et de ce que $\mathrm r^*$ est entier sur $\mathrm r$ que l'on a $\mathrm n\ne\mathrm q$. Comme $\mathrm n^*$ contient $\mathrm p^*=\mathrm q^*\cap\mathrm o^*$, il contient aussi $\mathrm q\cap\mathrm o=\mathrm p$, d'où $\mathrm p\subset\mathrm n$.

Supposons donc maintenant que $\mathrm o$ soit intégralement clos. On sait qu'il en est de même de $\mathrm o[t_1,..,t_n]$. L'idéal $\mathrm p\mathrm o[t_1,..,t_n]$ engendré par $\mathrm p$ dans $\mathrm o[t_1,..,t_n]$ est premier. Soit en effet $\mathrm R$
